% Figura: hmi-auto -- requiere % =====================================================================
%  Estilos compartidos por los diagramas TikZ de los capitulos.
%  Va en el PREAMBULO del documento principal:
%      % =====================================================================
%  Estilos compartidos por los diagramas TikZ de los capitulos.
%  Va en el PREAMBULO del documento principal:
%      % =====================================================================
%  Estilos compartidos por los diagramas TikZ de los capitulos.
%  Va en el PREAMBULO del documento principal:
%      \input{diagramas/estilos-diagramas}
% =====================================================================

\usepackage{tikz}
\usepackage{amsmath}
\usepackage{amssymb}
\usetikzlibrary{arrows.meta,positioning,shapes.geometric,shapes.misc,calc,fit,backgrounds,decorations.pathreplacing}

\definecolor{cbloque}{HTML}{E8EEF4}
\definecolor{cborde}{HTML}{4A6D8C}
\definecolor{cacento}{HTML}{D9E6D5}
\definecolor{cacentob}{HTML}{5A7A4A}
\definecolor{cgris}{HTML}{EDEDED}
\definecolor{cgrisb}{HTML}{888888}
\definecolor{calarma}{HTML}{F6DEDE}
\definecolor{calarmab}{HTML}{A03030}

\tikzset{
  bloque/.style={draw=cborde,fill=cbloque,rounded corners=2pt,align=center,
                 inner sep=5pt,font=\footnotesize,minimum height=9mm},
  bloqueg/.style={bloque,draw=cgrisb,fill=cgris},
  bloquea/.style={bloque,draw=cacentob,fill=cacento},
  bloquer/.style={bloque,draw=calarmab,fill=calarma},
  proceso/.style={bloque,minimum width=52mm},
  decision/.style={draw=cborde,fill=cbloque,align=center,font=\scriptsize,
                   diamond,aspect=2.2,inner sep=1pt},
  flecha/.style={-{Stealth[length=2mm]},draw=cborde!80!black,thick},
  flechag/.style={flecha,draw=cgrisb},
  et/.style={font=\scriptsize,inner sep=1.5pt,fill=white,align=center},
  etl/.style={et,midway},
  grupo/.style={draw=cgrisb,dashed,rounded corners=3pt,inner sep=4mm},
  tit/.style={font=\scriptsize\bfseries,text=cgrisb!60!black},
}

% --- utilidades para diagramas de secuencia ---
\tikzset{
  vida/.style={draw=cgrisb,dash pattern=on 2pt off 2pt},
  cab/.style={bloque,minimum width=26mm,font=\scriptsize\bfseries},
  nota/.style={draw=cacentob,fill=cacento,font=\scriptsize,align=center,
               rounded corners=1pt,inner sep=3pt},
}
\newcommand{\msg}[4]{\draw[flecha] (#1,#3) -- (#2,#3) node[et,midway,above=0pt]{#4};}
\newcommand{\msgr}[4]{\draw[flecha,dashed] (#1,#3) -- (#2,#3) node[et,midway,above=0pt]{#4};}
\newcommand{\auto}[3]{%
  \draw[flecha] (#1,#2) -- ++(0.55,0) -- ++(0,-0.42) -- (#1,{#2-0.42});
  \node[et,anchor=west] at ({#1+0.7},{#2-0.21}) {#3};}


% =====================================================================

\usepackage{tikz}
\usepackage{amsmath}
\usepackage{amssymb}
\usetikzlibrary{arrows.meta,positioning,shapes.geometric,shapes.misc,calc,fit,backgrounds,decorations.pathreplacing}

\definecolor{cbloque}{HTML}{E8EEF4}
\definecolor{cborde}{HTML}{4A6D8C}
\definecolor{cacento}{HTML}{D9E6D5}
\definecolor{cacentob}{HTML}{5A7A4A}
\definecolor{cgris}{HTML}{EDEDED}
\definecolor{cgrisb}{HTML}{888888}
\definecolor{calarma}{HTML}{F6DEDE}
\definecolor{calarmab}{HTML}{A03030}

\tikzset{
  bloque/.style={draw=cborde,fill=cbloque,rounded corners=2pt,align=center,
                 inner sep=5pt,font=\footnotesize,minimum height=9mm},
  bloqueg/.style={bloque,draw=cgrisb,fill=cgris},
  bloquea/.style={bloque,draw=cacentob,fill=cacento},
  bloquer/.style={bloque,draw=calarmab,fill=calarma},
  proceso/.style={bloque,minimum width=52mm},
  decision/.style={draw=cborde,fill=cbloque,align=center,font=\scriptsize,
                   diamond,aspect=2.2,inner sep=1pt},
  flecha/.style={-{Stealth[length=2mm]},draw=cborde!80!black,thick},
  flechag/.style={flecha,draw=cgrisb},
  et/.style={font=\scriptsize,inner sep=1.5pt,fill=white,align=center},
  etl/.style={et,midway},
  grupo/.style={draw=cgrisb,dashed,rounded corners=3pt,inner sep=4mm},
  tit/.style={font=\scriptsize\bfseries,text=cgrisb!60!black},
}

% --- utilidades para diagramas de secuencia ---
\tikzset{
  vida/.style={draw=cgrisb,dash pattern=on 2pt off 2pt},
  cab/.style={bloque,minimum width=26mm,font=\scriptsize\bfseries},
  nota/.style={draw=cacentob,fill=cacento,font=\scriptsize,align=center,
               rounded corners=1pt,inner sep=3pt},
}
\newcommand{\msg}[4]{\draw[flecha] (#1,#3) -- (#2,#3) node[et,midway,above=0pt]{#4};}
\newcommand{\msgr}[4]{\draw[flecha,dashed] (#1,#3) -- (#2,#3) node[et,midway,above=0pt]{#4};}
\newcommand{\auto}[3]{%
  \draw[flecha] (#1,#2) -- ++(0.55,0) -- ++(0,-0.42) -- (#1,{#2-0.42});
  \node[et,anchor=west] at ({#1+0.7},{#2-0.21}) {#3};}


% =====================================================================

\usepackage{tikz}
\usepackage{amsmath}
\usepackage{amssymb}
\usetikzlibrary{arrows.meta,positioning,shapes.geometric,shapes.misc,calc,fit,backgrounds,decorations.pathreplacing}

\definecolor{cbloque}{HTML}{E8EEF4}
\definecolor{cborde}{HTML}{4A6D8C}
\definecolor{cacento}{HTML}{D9E6D5}
\definecolor{cacentob}{HTML}{5A7A4A}
\definecolor{cgris}{HTML}{EDEDED}
\definecolor{cgrisb}{HTML}{888888}
\definecolor{calarma}{HTML}{F6DEDE}
\definecolor{calarmab}{HTML}{A03030}

\tikzset{
  bloque/.style={draw=cborde,fill=cbloque,rounded corners=2pt,align=center,
                 inner sep=5pt,font=\footnotesize,minimum height=9mm},
  bloqueg/.style={bloque,draw=cgrisb,fill=cgris},
  bloquea/.style={bloque,draw=cacentob,fill=cacento},
  bloquer/.style={bloque,draw=calarmab,fill=calarma},
  proceso/.style={bloque,minimum width=52mm},
  decision/.style={draw=cborde,fill=cbloque,align=center,font=\scriptsize,
                   diamond,aspect=2.2,inner sep=1pt},
  flecha/.style={-{Stealth[length=2mm]},draw=cborde!80!black,thick},
  flechag/.style={flecha,draw=cgrisb},
  et/.style={font=\scriptsize,inner sep=1.5pt,fill=white,align=center},
  etl/.style={et,midway},
  grupo/.style={draw=cgrisb,dashed,rounded corners=3pt,inner sep=4mm},
  tit/.style={font=\scriptsize\bfseries,text=cgrisb!60!black},
}

% --- utilidades para diagramas de secuencia ---
\tikzset{
  vida/.style={draw=cgrisb,dash pattern=on 2pt off 2pt},
  cab/.style={bloque,minimum width=26mm,font=\scriptsize\bfseries},
  nota/.style={draw=cacentob,fill=cacento,font=\scriptsize,align=center,
               rounded corners=1pt,inner sep=3pt},
}
\newcommand{\msg}[4]{\draw[flecha] (#1,#3) -- (#2,#3) node[et,midway,above=0pt]{#4};}
\newcommand{\msgr}[4]{\draw[flecha,dashed] (#1,#3) -- (#2,#3) node[et,midway,above=0pt]{#4};}
\newcommand{\auto}[3]{%
  \draw[flecha] (#1,#2) -- ++(0.55,0) -- ++(0,-0.42) -- (#1,{#2-0.42});
  \node[et,anchor=west] at ({#1+0.7},{#2-0.21}) {#3};}

 en el preambulo
\begin{tikzpicture}[x=1cm,y=1cm,
  conector/.style={draw=cborde,fill=cbloque,circle,minimum size=7mm,
                   inner sep=0pt,font=\footnotesize\bfseries}]

% ---------------- columna A: validación y puesta en el punto medio ----------------
\node[bloquea,text width=44mm] (a1) at (0,0)
  {\textbf{SINTONIZAR AUTO}\\[1pt]\scriptsize el operador presiona el botón};
\node[decision,text width=24mm] (a2) at (0,-1.8) {¿puerto\\abierto?};
\node[bloque,text width=48mm]  (a3) at (0,-3.6)
  {$t \leftarrow \dfrac{t_{\min P_{+}} + t_{\min P_{-}}}{2}$};
\node[decision,text width=26mm] (a4) at (0,-5.5) {¿hay mínimos\\registrados?};
\node[decision,text width=26mm] (a5) at (0,-7.5) {$t\in[h_{\min},h_{\max}]$?};
\node[bloque,text width=48mm]  (a6) at (0,-9.4)
  {Enviar \texttt{s}\\[1pt]\scriptsize el heater deja de barrer};
\node[bloque,text width=48mm]  (a7) at (0,-10.9)
  {Enviar \texttt{T<t>} y \texttt{w}\\[1pt]\scriptsize el setpoint entra sin esperar al ciclo en curso};
\node[bloque,text width=48mm]  (a8) at (0,-12.5)
  {Esperar \texttt{AUTO\_SETTLE\_TIME\_S} $=20$\,s\\[1pt]\scriptsize estabilización térmica};
\node[conector] (ca) at (0,-14.0) {A};

% salidas por rechazo
\node[bloquer,text width=40mm] (f1) at (-7.2,-1.8)
  {\scriptsize Mensaje \emph{No conectado} en el monitor serie};
\node[bloquer,text width=40mm] (f2) at (-7.2,-5.5)
  {\scriptsize Estado \textbf{SIN MÍNIMOS}: barrer con AVANZAR o RETROCEDER primero};
\node[bloquer,text width=40mm] (f3) at (-7.2,-7.5)
  {\scriptsize Estado \textbf{FUERA DE RANGO}: la temperatura media cae fuera de los límites};

\draw[flecha] (a1) -- (a2);
\draw[flecha] (a2) -- node[et,right]{sí} (a3);
\draw[flecha] (a3) -- (a4);
\draw[flecha] (a4) -- node[et,right]{sí} (a5);
\draw[flecha] (a5) -- node[et,right]{sí} (a6);
\draw[flecha] (a6) -- (a7);
\draw[flecha] (a7) -- (a8);
\draw[flecha] (a8) -- (ca);
\draw[flecha] (a2) -- node[et,above]{no} (f1);
\draw[flecha] (a4) -- node[et,above]{no} (f2);
\draw[flecha] (a5) -- node[et,above]{no} (f3);

% ---------------- columna B: medición y enganche ----------------
\node[conector] (cb) at (9.2,-0.7) {A};
\node[bloque,text width=48mm] (b1) at (9.2,-2.2)
  {\texttt{auto\_cycles\_left} $\leftarrow$ \texttt{AUTO\_MEASURE\_CYCLES} $=4$};
\node[bloque,text width=48mm] (b2) at (9.2,-4.0)
  {Esperar el cierre de un ciclo\\[1pt]\scriptsize línea \texttt{[log]} del firmware};
\node[bloque,text width=48mm] (b3) at (9.2,-5.6)
  {\texttt{auto\_cycles\_left} $\leftarrow$ \texttt{auto\_cycles\_left} $-\,1$};
\node[decision,text width=24mm] (b4) at (9.2,-7.5) {\texttt{auto\_cycles}\\\texttt{\_left} $>0$?};
\node[bloque,text width=54mm] (b5) at (9.2,-9.6)
  {Enviar \texttt{g}\\[1pt]\scriptsize el lock toma como referencia el $\mathit{prt}$ de dos ciclos atrás, ya estabilizado};
\node[bloquea,text width=44mm] (b6) at (9.2,-11.5)
  {\textbf{SINTONIZADO}};

\draw[flecha] (cb) -- (b1);
\draw[flecha] (b1) -- (b2);
\draw[flecha] (b2) -- (b3);
\draw[flecha] (b3) -- (b4);
\draw[flecha] (b4) -- node[et,right]{no} (b5);
\draw[flecha] (b5) -- (b6);
\draw[flecha] (b4.east) -- node[et,pos=0.28,above]{sí} (13.4,-7.5) -- (13.4,-4.0) -- (b2.east);

% ---------------- notas ----------------
\node[nota,text width=58mm,anchor=north west] at (-9.7,2.5)
 {\textbf{Precondición.} Durante el barrido con \texttt{f} o \texttt{b} el HMI sigue
  ciclo a ciclo los mínimos de $p^{+}$ y $p^{-}$ y la temperatura a la que ocurren.
  DETENER (\texttt{s}) los borra.};

\node[nota,text width=62mm,anchor=north west] at (3.6,-12.6)
 {\textbf{Cancelación.} En cualquier punto, un comando de \texttt{INTERRUPTING\_COMMANDS}
  (\texttt{s}, \texttt{f}, \texttt{b}, \texttt{g}, \texttt{e}, \texttt{T}), un botón de
  modo o la caída del puerto cancelan la secuencia.};

\end{tikzpicture}
